\rulesection{12.0 Combat}

During the Combat Phase, combat may occur in any hex occupied by both sides. The first player determines the order in which these battles are resolved.

Compare initiative of each side's senior leader. The player with the higher initiative may decide whether to attack first. If initiatives are equal, the player moving into the hex second decides whether to attack first. If the player with first choice declines, his opponent may attack. If neither player attacks, there is no combat and all units remain in the hex. If there is combat, use the procedure below unless neither side has more than 10 strength points or both sides have only cavalry; in that case use skirmish combat.

\rulesubsection{12.1 Procedure}

\begin{itemize}
\item Each player checks supply.
\item Each player chooses a combat option and checks for wagon depletion.
\item Players compare combat options and determine the number of rounds to be fought.
\item Each player determines troop commitment.
\item Rounds are fought.
\item Retreats are executed. Each player may reorient facings of involved cavalry units.
\item In some scenarios, the winner determines whether a major or minor victory was achieved.
\end{itemize}

\rulesubsection{12.2 Participating Units}

If a player attacks, he may use some or all of his units, except that demoralized units may never attack. Units that do not attack are not affected by the battle. All defending units always take part.

\rulesubsection{12.3 Combat Supply}

A combat unit is in supply if an undepleted wagon is within 2 MPs, a supply base is within 4 MPs, or a depot is within 4 MPs. A single wagon, base, or depot may supply any number of units and any number of combats in a single Combat Phase.

\textbf{12.31 Tracing Supply Lines.} Trace from the unit to the supply source. Count terrain costs using the combat unit's terrain and ZOC costs. Railroads count as roads, but rail movement may not be used for supply. A supply line may not include a hex occupied by an enemy combat unit other than the hex of the units being supplied.

\textbf{12.32 Effects.} A combat unit unable to trace a supply line of the correct length is unsupplied.

\textbf{12.321} If a player's units are unsupplied, he may not choose the \textit{Pitched Battle} or \textit{Assault} combat options.

\textbf{12.322} Unsupplied units must retreat at the end of the final combat round unless the enemy is forced to retreat because of demoralization.

\textbf{12.33 Wagon Depletion.} A wagon used to supply combat may become depleted. Immediately after choosing a combat option, roll one die. If the option is \textit{Pitched Battle}, a roll of 5 or 6 depletes the wagon. If any other option is chosen, only a roll of 6 depletes it. Place a Depletion marker on a wagon that becomes depleted. At the end of the Combat Phase the wagon is flipped to its depleted side and the marker removed. A player may not voluntarily choose to be unsupplied if a supply line is possible.

\rulesubsection{12.4 Combat Options}

In resolving each combat, both attacker and defender choose an option chit secretly and reveal after both have chosen.

\begin{itemize}
\item Attacker options: \textit{Probe}, \textit{Assault}, \textit{Pitched Battle}
\item Defender options: \textit{Withdraw}, \textit{Stand}, \textit{Pitched Battle}
\end{itemize}

\textbf{12.41 Restrictions.} If a player's units are unsupplied, he may not choose \textit{Pitched Battle} or \textit{Assault}. If his units have no legal retreat path, he may not choose \textit{Withdraw}.

\textbf{12.42 Effects.} Cross-reference the two choices on the Combat Options Chart. Results of \texttt{1} or \texttt{2} mean the combat lasts that many rounds. Result \texttt{U} means rounds continue until one or both sides become demoralized.

\textbf{12.421} \textit{Pitched Battle} increases the player's Commitment Table die roll by 1 and also makes wagon depletion more likely.

\textbf{12.422} A player who chooses \textit{Withdraw} must retreat at the end of the battle and reduces his Commitment Table die roll by 1.

\textbf{12.423} \textit{Probe} reduces the player's Commitment Table die roll by 1.

\rulesubsection{12.5 Commitment}

Immediately before combat rounds begin, each player determines commitment. Select the senior friendly leader in the hex, modify his command rating if necessary, select the corresponding column on the Commitment Table, roll two dice, add any modifiers from the table, and cross-reference the result. The result is the percentage of the force's total SPs that may be used in the combat.

Commitment occurs only once, before rounds begin, and represents the arrival of fresh troops throughout the battle as well as those initially sent into action.

\rulesubsection{12.6 Resolving Rounds}

To resolve a round, each player totals and modifies the strength of his committed units, locates the corresponding Combat Results Table column, and shifts one column left if his opponent is defending in an entrenchment. Each player then rolls one die, modifies the result as listed with the table, and cross-references with the proper column. The result is the number of enemy SPs eliminated. The enemy player chooses which of his units suffer losses.

\textbf{12.61} Although it is easier to resolve fire in sequence, each round's results take effect simultaneously.

\textbf{12.62 Strength Modifications.}
\begin{itemize}
\item Halve a unit if it entered the hex during the current turn through an adjacent bridge or river hexside.
\item Halve a unit if combat occurs in a ridge hex and it entered during the current turn through a slope hexside.
\item Double a unit if it is defending and entrenched.
\end{itemize}

Modifications are cumulative. Apply them to the total strength of all units to which they apply, rounding to the nearest whole number and rounding one-half up.

\textbf{12.63 Die-Roll Modifiers.}
\begin{itemize}
\item Add the combat bonus of the senior leader on the firing side.
\item Add 1 if the player chose \textit{Pitched Battle}.
\item Subtract 1 if the player chose \textit{Withdrawal} or \textit{Probe}.
\item Add 1 to all fire by the defender in all battles.
\end{itemize}

\textbf{12.64 Excess Strength.} If committed SPs exceed 56, roll once using the 56 column and again for the excess. A player with more than 56 committed SPs must use the 56 column.

\textbf{12.65 Excess Losses.} If losses exceed committed SPs, some of the excess may be removed from uncommitted SPs. The other player rolls a die, adds the number of excess losses, and consults the Excess Losses Table to determine additional losses.

\rulesubsection{12.7 Skirmish Combat}

When both sides in a combat are cavalry or 10 SPs or less of infantry, use skirmish combat:
\begin{itemize}
\item Each player totals his force's SPs.
\item Each player rolls on the Combat Results Table and applies the result.
\item Each side that lost at least one SP checks for demoralization. Units that become demoralized must retreat.
\item Units without supply must retreat. Wagons are never depleted due to skirmish combat.
\item Combat is completed after a single round.
\end{itemize}
